\documentclass[11pt]{article}
\usepackage[utf8]{inputenc}
\usepackage[T1]{fontenc}
\usepackage{lmodern}
\usepackage[margin=0.9in]{geometry}
\usepackage{enumitem}
\usepackage{xcolor}
\usepackage{hyperref}
\usepackage{amsmath, amssymb}
\usepackage{booktabs}
\usepackage[normalem]{ulem}

\hypersetup{
  colorlinks = true,
  linkcolor = blue!50!black,
  urlcolor = blue!50!black
}

\definecolor{keyblue}{RGB}{19,41,75}      % UNC navy --- main emphasis
\definecolor{keycar}{RGB}{75,156,211}     % Carolina blue --- headline numbers

\setlist[itemize]{leftmargin=1.4em, itemsep=0.25em, topsep=0.35em}
\setlength{\parindent}{0pt}
\setlength{\parskip}{0.55em}
\newcommand{\scriptlabel}{\textbf{Script. }}
\newcommand{\sectionlabel}[1]{\textbf{#1}}
\newcommand{\actioncue}[1]{\textit{\textbf{#1}}}
% Emphasis macros for use inside scripts:
\newcommand{\key}[1]{\textcolor{keyblue}{\textbf{#1}}}        % bold + navy: say-this points
\newcommand{\num}[1]{\textcolor{keycar}{\textbf{\uline{#1}}}} % bold + underline + Carolina: headline numbers
\title{Supervised Bayesian Matrix Factorization\\Identifies Prognostic Gene Expression Programs\\[0.8em]\large Presentation Notes}
\author{Andrew Walther}
\date{June 18, 2026}

\begin{document}
\maketitle

\noindent This companion is written as a preparation guide for an approximately 25-minute lab
meeting talk. The main-deck sections are written as the primary script. The figure-and-table
analysis guide, the quick review prompts, and the anticipated-questions section are reference
material to study before the talk and to draw on during discussion. The appendix scripts are short
and are intended only for slides that come up if the discussion turns technical.

\section*{Executive Summary}

Pancreatic ductal adenocarcinoma (PDAC) needs validated molecular prognosis, but the conventional
\emph{two-step} pipeline --- run an unsupervised factorization (PCA / EBMF), then test factors for
survival after the fact --- chases the largest expression variance and misses low-variance programs
that cleanly separate survivors. This work develops a \emph{joint, supervised Bayesian matrix
factorization}: a single set of patient program activations both reconstructs the expression matrix
and drives a Cox risk score, so survival shapes the factors \emph{during} estimation. New patients
are scored by an exact genomic projection ($\eta = (Y_{\text{new}}F)\beta$), so out-of-sample
scoring is leakage-free with no re-fitting, and the model is fully Bayesian with adaptive priors and
no genomics-versus-survival weight to tune.

\begin{itemize}
\item \sectionlabel{Findings:} On a controlled simulation, the supervised model reaches a held-out
C-index of $0.81$ in the realistic scenario versus $0.72$ for the two-step, while all methods sit at
chance ($\approx 0.54$) in a null scenario --- no fabricated signal. On real PDAC data (train
TCGA-PAAD + CPTAC, validate five fully held-out cohorts), it attains mean external C $= 0.636$,
beating the two-step in all five cohorts.
\item \sectionlabel{Parsimony:} Of seven retained programs, only two carry survival signal --- an
adverse MET/integrin/glycolytic program and a protective epithelial-differentiation program.
\item \sectionlabel{Impact:} Supervision recovers prognostic structure the two-step misses, on both
synthetic and independent real cohorts; preprocessing (per-platform normalization before merging) is
decisive for keeping the survival coefficients from collapsing; and the output is two interpretable
programs with exact external scoring. Next step: pathway enrichment to assign biology, en route to a
methods-primary manuscript.
\end{itemize}

\section*{Opening}

\noindent\emph{Slide numbers below match the deck's on-screen numbering. The three section-divider
slides (3 Background, 5 Methods, 13 Results) count toward that numbering; appendix slides are
``uncounted'' in the deck and so are labeled A1--A3 here.}

\noindent\textbf{Pacing.} 20--25 minutes over 24 slides is about \key{one minute per slide}. Each
\scriptlabel paragraph is written to be read in 30--45 seconds --- keep moving. Text in
\key{navy bold} is a must-say point; \num{underlined Carolina blue} marks the headline numbers.

\subsection*{Slide 1 --- Title Slide}
\scriptlabel Thanks everyone --- a project update on the supervised Bayesian matrix factorization
work. The one-line version: \key{bringing survival into the factorization changes which gene
programs the model prioritizes}, and it \key{beats the unsupervised two-step} both on simulation and
across \num{five} independent PDAC cohorts.

\begin{itemize}
\item \sectionlabel{Main Takeaway:} Frame the talk as a methods update \emph{and} a validated
results update --- this is no longer preliminary.
\item \sectionlabel{Emphasis:} The talk is about clinically relevant latent structure, not
reconstruction quality. Set that expectation early.
\item \sectionlabel{Timing:} Two sentences, then move into the agenda.
\end{itemize}

\section*{Main Deck}

\subsection*{Slide 2 --- Agenda}
\scriptlabel A couple of minutes on motivation, then the method, then \key{most of the talk is
results} --- a multi-cohort simulation, the real-data gene programs and survival, and external
validation. I close with impact and next steps.

\begin{itemize}
\item \sectionlabel{Main Takeaway:} This is a results-forward talk; the method section directly
answers the questions raised in April.
\item \sectionlabel{Emphasis:} Flag that the heart of the talk is the synthetic-versus-real
comparison against the unsupervised two-step.
\end{itemize}

\subsection*{Slide 3 --- Section divider: Background}
\scriptlabel (Divider slide --- no content.) Pause here to transition from agenda into motivation.

\subsection*{Slide 4 --- Background: PDAC needs molecular prognosis}
\scriptlabel PDAC has five-year survival under 12\% and \key{no adopted molecular stratification at
diagnosis}. The signal is there but mixed across malignant, stromal, and immune compartments and
often low-variance. Our goal: recover \key{gene expression programs} --- a column of $F$ with
patient activation in $L$ --- that predict survival, validated on independent cohorts. Why
\key{jointly} rather than two-step: unsupervised axes chase the largest variance, often batch, so a
\key{15\%-variance program may be batch while a 3\% one separates survivors}. Supervision steers the
factors toward survival-relevant structure during estimation.

\begin{itemize}
\item \sectionlabel{Main Takeaway:} There is a genuine clinical need, and the methodological hook
--- joint over two-step --- is the thread of the entire talk.
\item \sectionlabel{Figure/Table Purpose:} The callout contrasts ``describing tumor heterogeneity''
with ``predicting disease course.''
\item \sectionlabel{Good Line to Stress:} ``A high-variance program can be batch effect; a
low-variance program can be the one that separates survivors. Supervision prioritizes the second.''
\end{itemize}

\subsection*{Slide 5 --- Section divider: Methods}
\scriptlabel (Divider slide --- no content.) Signal the shift from motivation into the method.

\subsection*{Slide 6 --- The model --- two coupled components}
\scriptlabel Two coupled pieces: a low-rank factorization $Y = LF^\top + E$ for expression, and a
Cox model $h_i(t) = h_0(t)\exp(\eta_i)$ for survival, driven by the \key{same patient loadings $L$}.
The table defines every symbol. The key point: a program is \key{prognostic only if $\beta_k \neq
0$} --- otherwise it explains expression alone. And the \key{baseline hazard $h_0(t)$ is
non-parametric}: it cancels in the partial likelihood, so no Weibull/exponential assumption, just
proportional hazards.

\begin{itemize}
\item \sectionlabel{Main Takeaway:} One latent space drives both genomics reconstruction and the
survival score.
\item \sectionlabel{Figure/Table Purpose:} The dimension table defines every object once so later
slides move faster.
\item \sectionlabel{Emphasis:} Two specific April requests are answered here --- every parameter is
defined, and $h_0(t)$ is non-parametric and cancels.
\item \sectionlabel{Good Line to Stress:} ``$\beta_k = 0$ means the model keeps that program for
expression but not for prognosis --- that separation is the whole point.''
\end{itemize}

\subsection*{Slide 7 --- Two linear predictors --- $L\beta$ vs $(YF)\beta$}
\scriptlabel The key update since April: \key{two linear predictors}. The loadings predictor
$\eta = L\beta$ is fully joint, but scoring a new patient needs an \key{OLS re-projection --- an
approximation}. The projection predictor $\eta = (YF)\beta$ scores new patients by the \key{identical
training formula --- exact, no mismatch}. On the priors --- now stated correctly: $q(\cdot)$ is the
variational posterior for each block, fit by empirical Bayes every update. \key{Point-exponential
(non-negative) on $L$ and $F$, normal on $\beta$}, closed-form MLE for $\tau$ --- \key{not
point-normal}.

\begin{itemize}
\item \sectionlabel{Main Takeaway:} Report both predictors; the projection $(YF)\beta$ is the
recommended one because it gives exact, leakage-free external scoring.
\item \sectionlabel{Emphasis:} The priors are \emph{point-exponential} on $L,F$ and \emph{normal}
on $\beta$. This is a correction from April --- they are \emph{not} point-normal. Say this clearly.
\item \sectionlabel{Good Line to Stress:} ``$(YF)\beta$ scores a held-out patient with the same
formula we trained on --- there is no projection approximation to worry about.''
\item \sectionlabel{Likely Question:} ``Why keep $L\beta$ at all?'' It is the fully-joint reference:
survival shapes the loadings most directly there, and comparing the two quantifies what we give up
for exact scoring (in practice, very little --- see the external C-index slide).
\end{itemize}

\subsection*{Slide 8 --- The objective --- ELBO \& variational family}
\scriptlabel The variational family is \key{mean-field} --- every block factorizes independently.
Inference maximizes the \key{ELBO}, with three pieces: a Gaussian genomics-fit term, a survival-fit
term (a \key{2nd-order Taylor expansion of the Cox likelihood} that keeps it conjugate to the
updates), and a per-block KL term that is the adaptive shrinkage. The \key{shared $L$ enters both
likelihood terms --- that coupling is the supervision}. And it balances the two automatically:
\key{no $\alpha$ to tune} ($\lambda = 1$).

\begin{itemize}
\item \sectionlabel{Main Takeaway:} Supervision is built into a single joint objective; the shared
$L$ in both likelihood terms is the mechanism.
\item \sectionlabel{Emphasis:} State plainly what the ELBO is \emph{for}: it is the fitting
objective that CAVI maximizes, \emph{not} the criterion used to choose the number of programs (that
is cross-validation, two slides later).
\item \sectionlabel{Good Line to Stress:} ``Both likelihoods share the same latent loadings, so the
factors cannot specialize to only one task.''
\item \sectionlabel{Likely Question:} ``Why is there no $\alpha$ to weight genomics against
survival?'' Because the ELBO is a single coherent objective on the same probability scale; the two
likelihood terms and the KL penalty are already commensurate, so $\lambda = 1$. This was a specific
April concern about hand-tuning the balance.
\end{itemize}

\subsection*{Slide 9 --- Inference --- factor-wise CAVI}
\scriptlabel Inference is \key{coordinate-ascent variational inference}: one factor at a time, each
update an \key{EBNM sub-problem}. The one detail to land: the loadings precision is
\key{dual-source} --- a genomics piece plus a \key{survival piece $W_{ii}E[\beta_k^2]$} --- so
survival genuinely shapes the loadings. Convergence, in plain English: stop when the \key{relative
ELBO change falls below $10^{-5}$} ($\Delta L$, $\Delta\beta$ are diagnostics only). And a
\key{sign correction}: concordance is sign-invariant, so we orient each program by \key{training}
concordance and reuse it at test.

\begin{itemize}
\item \sectionlabel{Main Takeaway:} The algorithm updates one factor at a time, each block reducing
to an EBNM call; convergence is judged on the ELBO itself.
\item \sectionlabel{Emphasis:} Convergence is the \emph{relative ELBO change below $10^{-5}$} --- the
ELBO is both the objective and the convergence monitor. The $\Delta L$ / $\Delta\beta$ thresholds are
diagnostics, not the stopping rule.
\item \sectionlabel{Good Line to Stress:} ``The loadings precision literally adds a survival term to
the genomics term --- that is where supervision enters the updates.''
\item \sectionlabel{Likely Question:} ``Is the convergence rule about the \emph{mean} of $L$ and
$\beta$?'' No --- the stopping rule is the relative change in the ELBO. We additionally watch the
mean absolute per-element change in $L$ and $\beta$ as a sanity check; averaging avoids a single
oscillating factor blocking the stop, but it is diagnostic only.
\end{itemize}

\subsection*{Slide 10 --- Preprocessing is crucial --- per-platform normalization}
\scriptlabel \key{The most important practical finding since April.} Merging platforms leaves genes
on different scales. The survival precision is $A_k = \sum_i w_i (YF)_{ik}^2$ --- $w_i$ the Cox
working weight, $(YF)_{ik}$ the projection score. Without correction, \key{between-platform variance
swamps $A_k$ and the survival coefficients collapse to $\beta \to 0$}. The fix: \key{z-standardize
each platform \emph{before} merging} (with log2 + per-cohort gene selection, leaving \num{2064}
genes). In an 18-config benchmark, \key{10 of 12} non-per-platform pipelines collapse. And no
leakage --- normalization is fit within each cohort before any split.

\begin{itemize}
\item \sectionlabel{Main Takeaway:} Per-platform normalization before merging is the difference
between a working model and $\beta \to 0$.
\item \sectionlabel{Emphasis:} Define both pieces of $A_k$ aloud --- $w_i$ is the Cox working weight,
$(YF)_{ik}$ is the projection score. The audience flagged this expression as under-explained.
\item \sectionlabel{Good Line to Stress:} ``Ten of twelve alternative pipelines collapse the
survival signal to zero --- this preprocessing step is load-bearing, not cosmetic.''
\item \sectionlabel{Likely Question:} ``Isn't per-cohort normalization leakage?'' No --- it is fit
within each training cohort before any train/test split, and each held-out cohort is normalized by
its own statistics at scoring time, never using training or test labels.
\end{itemize}

\subsection*{Slide 11 --- Selecting the number of programs --- $K$ and $K_{\text{eff}}$}
\scriptlabel Two distinct numbers, kept separate. \key{$K$ = how many programs we retain}, chosen by
\key{5-fold cross-validation} (1-SE rule, biological floor $K \ge 3$) $\to$ \num{$K = 7$}.
\key{$K_{\text{eff}}$ = how many are prognostic}: a program is survival-active if \key{$|\hat\beta_k|
> 10^{-3}$}, and only \num{$K_{\text{eff}} = 2$} of the seven clear it --- the rest explain
expression only. \key{$K$ is reconstruction capacity; $K_{\text{eff}}$ is the prognostic subset} ---
always quote both.

\begin{itemize}
\item \sectionlabel{Main Takeaway:} $K = 7$ by cross-validation; $K_{\text{eff}} = 2$ by the
survival-coefficient threshold. Always quote both.
\item \sectionlabel{Emphasis:} $K$ is now chosen by CV --- this directly answers April's ``how is
$K$ selected?'' The two filters are different criteria, not the same number.
\item \sectionlabel{Good Line to Stress:} ``$K$ is how many programs the model keeps;
$K_{\text{eff}}$ is how many of those actually predict survival.''
\item \sectionlabel{Likely Question:} ``Why a floor of three?'' It is a biological/comparison
minimum aligned with the lab's penalized survival-MF setup --- we want enough programs to compare
shared and specific structure, not to discover the model wants only one.
\end{itemize}

\subsection*{Slide 12 --- $L$ extension --- cohort-membership indicator}
\scriptlabel An extension we evaluated and \key{do not recommend}. We append dummy columns to $L$
(one per non-reference cohort) with \key{$\beta$ fixed at zero}, so they absorb platform offsets
while survival uses only the biological programs. The forced dummy costs a degree of freedom, so
\key{$K: 7 \to 8$}. But under per-platform z-standardization the \key{offset is already removed} ---
the indicator just spends a slot, moving external C \key{\emph{down} \num{0.636} $\to$ \num{0.614}}.
So: \key{no cohort indicator}.

\begin{itemize}
\item \sectionlabel{Main Takeaway:} Recommendation is \emph{no} cohort indicator under per-platform
z-standardization --- it costs a factor and slightly lowers external C.
\item \sectionlabel{Emphasis:} Note explicitly that the indicator requires one extra factor (a
degree of freedom), so $K: 7 \to 8$, to preserve the biological program count.
\item \sectionlabel{Good Line to Stress:} ``The indicator and the normalization solve the same
problem; once you normalize per platform, the indicator is redundant.''
\item \sectionlabel{Likely Question:} ``So when \emph{would} you use it?'' Only if you could not
normalize per platform --- e.g.\ a single merged matrix with unknown batch structure. Here we can,
so we do not need it.
\end{itemize}

\subsection*{Slide 13 --- Section divider: Results}
\scriptlabel (Divider slide --- no content.) Mark the hand-off from method into results.

\subsection*{Slide 14 --- Data --- synthetic validation \& PDAC cohorts}
\scriptlabel Two datasets. The simulation is a \key{completed} multi-cohort study with known ground
truth --- shared factors (carrying the survival signal) plus study-specific factors --- across
\key{three scenarios} (all-shared, an \key{all-cohort-specific negative control}, and the realistic
mix) and five model arms including the unsupervised EBMF baseline. On real data we \key{train on
TCGA-PAAD + CPTAC} ($n \approx 273$) and \key{validate on five fully held-out cohorts} scored by the
exact projection formula, spanning microarray and RNA-seq.

\begin{itemize}
\item \sectionlabel{Main Takeaway:} The validation is cross-cohort and cross-platform on both
simulated and real data --- not a single-dataset demonstration.
\item \sectionlabel{Figure/Table Purpose:} The cohort table grounds the audience in which cohorts
are held out and on what platform.
\item \sectionlabel{Emphasis:} The simulation is \emph{completed} (drop any ``in progress''
language), and the all-cohort-specific scenario is a genuine negative control.
\item \sectionlabel{Likely Question:} ``Why these five external cohorts?'' They are the publicly
available PDAC cohorts with matched expression and survival and sufficient sample size, and they
span microarray and RNA-seq so the test is platform-agnostic by construction.
\end{itemize}

\subsection*{Slide 15 --- Synthetic results --- shared \& cohort-specific factor recovery}
\scriptlabel Two complementary questions. \key{Figure 1 --- recovery}: how faithfully each program is
reconstructed (a correlation). \key{Figure 2 --- specificity accuracy}: whether the model correctly
\emph{labels} each program shared vs.\ study-specific (a classification). We need \key{both} ---
good programs with the wrong architecture would show as high recovery, low typing. The projection
model is \key{high on both} (recovery \num{0.90}--\num{0.92}, typing up to \num{1.00}), which is why
the cohort indicator adds little.

\begin{itemize}
\item \sectionlabel{Main Takeaway:} The projection model both reconstructs and correctly types the
shared prognostic programs, with no indicator columns.
\item \sectionlabel{Figure/Table Purpose:} Figure 1 = reconstruction fidelity (a correlation);
Figure 2 = structural classification (an accuracy). They measure different things.
\item \sectionlabel{Emphasis:} Motivate why \emph{both} plots are needed --- good recovery does not
guarantee correct shared-vs-specific typing.
\item \sectionlabel{Likely Question:} ``What is specificity-classification accuracy exactly?'' It is
the fraction of recovered factors correctly called shared vs.\ study-specific, judged from how a
factor's loading energy is distributed across cohorts --- shared factors load in all cohorts,
specific factors in one.
\end{itemize}

\subsection*{Slide 16 --- Synthetic results --- factor recovery heatmap}
\scriptlabel This makes recovery concrete. \key{Left: the true loadings}; \key{right: the estimate}
from the projection model with \key{no cohort indicator}. \key{Shared factors light up in both
cohort blocks, study-specific in one} --- and the estimate reproduces that structure straight from
the data. This is the visual basis for the previous slide's numbers.

\begin{itemize}
\item \sectionlabel{Main Takeaway:} The model recovers the intended block-sparse shared/specific
structure straight from the data, without being told which factors are shared.
\item \sectionlabel{Figure/Table Purpose:} Side-by-side true vs.\ estimated loadings; the visual
counterpart to the correlation and accuracy bars.
\item \sectionlabel{Good Line to Stress:} ``No indicator columns --- the model figures out shared
versus cohort-specific on its own.''
\end{itemize}

\subsection*{Slide 17 --- Synthetic results --- supervision outperforms the unsupervised two-step}
\scriptlabel \key{The headline of the simulation.} In the realistic scenario (mean over 5 seeds),
the supervised projection model reaches \num{0.81}, ahead of loadings at \num{0.76} and the
\key{unsupervised two-step at \num{0.72}}. And in the \key{negative control, every arm sits at
$\approx$\num{0.54}} --- \key{no false signal manufactured when none exists}. Clean, controlled
evidence that supervision helps and does not hallucinate.

\begin{itemize}
\item \sectionlabel{Main Takeaway:} Joint supervision beats the unsupervised two-step on data with
known ground truth, and invents no signal when there is none.
\item \sectionlabel{Figure/Table Purpose:} The C-index bars now \emph{include} the EBMF two-step, so
the comparison is on one plot.
\item \sectionlabel{Emphasis:} Say ``hybrid (shared \& cohort-specific)'' and ``mean over five
seeds'' precisely. The $0.54$ control is as important as the $0.81$.
\item \sectionlabel{Good Line to Stress:} ``Supervision buys about nine C-index points over the
two-step in the realistic scenario --- and nothing in the null, which is exactly what you want.''
\end{itemize}

\subsection*{Slide 18 --- PDAC cohorts --- prognostic gene programs}
\scriptlabel On real data, the recommended model retains \num{7} programs but only \num{2} are
survival-active (starred): \key{Program 3 protective, Program 7 adverse}. Weights are non-negative
(point-exponential prior). The point to land: the \key{largest expression program --- a stromal
collagen block --- is \emph{not} survival-active}. \key{Supervision separates expression-only from
prognostic structure} --- exactly what PCA-then-test cannot do.

\begin{itemize}
\item \sectionlabel{Main Takeaway:} Seven programs retained, two survival-active (Program 3
protective, Program 7 adverse); the biggest expression program is not prognostic.
\item \sectionlabel{Figure/Table Purpose:} The starred heatmap shows which programs are prognostic
and that they are non-negative, interpretable gene sets.
\item \sectionlabel{Emphasis:} State the model by its characteristics, not a code name. Quote
$K = 7$ retained, $K_{\text{eff}} = 2$ active.
\item \sectionlabel{Likely Question:} ``You said two active --- the heatmap looks like more are
starred?'' Only two clear the $|\hat\beta| > 10^{-3}$ threshold; any others are below it and
expression-only. Program 3 (protective) and Program 7 (adverse) are the two.
\end{itemize}

\subsection*{Slide 19 --- PDAC cohorts --- prognostic association}
\scriptlabel Both active programs \key{stratify patients into clear survival groups}, split at the
\key{median projection score $(YF)_k$} --- the quantity the model scores on. \key{Program 7
(adverse)} carries \key{MET, ITGA3, and glycolytic enzymes} --- an aggressive signature --- and does
worse, \key{$p = \num{0.00047}$}. \key{Program 3 (protective)} carries epithelial-differentiation
markers, $p = \num{0.0077}$. Directions read off the curves; \key{PH not violated} (Schoenfeld).

\begin{itemize}
\item \sectionlabel{Main Takeaway:} Both active programs produce clean, significant survival
separation on the score the model uses.
\item \sectionlabel{Figure/Table Purpose:} Kaplan--Meier curves, split at the median $(YF)_k$
projection score, with log-rank $p$-values and the empirical direction labeled.
\item \sectionlabel{Emphasis:} Name a few genes under each program --- MET/ITGA3/glycolytic for
adverse, epithelial-differentiation for protective --- to set up the pathway-enrichment next step.
\item \sectionlabel{Likely Question (be ready):} ``Does the joint $\hat\beta$ sign agree with the
adverse direction?'' For Program 7 the joint-model coefficient sign can run \emph{opposite} to its
marginal adverse direction --- a suppression effect among correlated programs. We therefore label
each program by the empirical survival direction the KM curve shows, and stratify by the $(YF)$
projection, not the raw $\beta$ sign. Be precise here; this is a subtle point that can otherwise look
like an inconsistency.
\end{itemize}

\subsection*{Slide 20 --- PDAC cohorts --- external C-index}
\scriptlabel The \key{real-data analog of the synthetic headline}, scored on five held-out cohorts.
The \key{projection model averages \num{0.636}}, loadings \num{0.622}, the \key{unsupervised
two-step \num{0.564}}. The projection model \key{wins in \emph{all five} cohorts}. Beyond the number,
I lean on \key{parsimony}: two active programs and exact external scoring, versus $\sim$20
unsupervised factors.

\begin{itemize}
\item \sectionlabel{Main Takeaway:} The projection supervised model beats the unsupervised two-step
across all five external cohorts; mean external C $= 0.636$.
\item \sectionlabel{Figure/Table Purpose:} Per-cohort + mean bars for the three arms, plotted at a
proper aspect ratio with the chance line.
\item \sectionlabel{Emphasis:} Use the labels ``projection (YFB) model'' and ``loadings (LB)
model.'' State the all-five-cohort win, not a vague ``competitive.''
\item \sectionlabel{Good Line to Stress:} ``The projection-based supervised joint model outperforms
the unsupervised two-step across all five external PDAC cohorts.''
\end{itemize}

\subsection*{Slide 21 --- Model comparison summary}
\scriptlabel One table, every model by its characteristics. The \key{recommendation: projection
predictor, per-platform z-std, no cohort indicator} --- \key{best on both synthetic (\num{0.81}) and
real external (\num{0.636})}, the most parsimonious structure ($K = 7 \to K_{\text{eff}} = 2$), exact
scoring. Loadings is close (\num{0.76} / \num{0.622}); the indicator costs a factor for no gain; the
two-step trails (\num{0.72} / \num{0.564}). The gene-selection recipe is \key{our own work},
motivated by the lab's penalized survival-MF.

\begin{itemize}
\item \sectionlabel{Main Takeaway:} The recommended configuration wins on both axes with the
simplest, most interpretable structure.
\item \sectionlabel{Figure/Table Purpose:} One table comparing linear predictor, preprocessing,
cohort indicator, $K$, $K_{\text{eff}}$, synthetic C, and real external C across the four models.
\item \sectionlabel{Emphasis:} Describe the gene-selection step as our own work; mention the
penalized-MF motivation only as a side note, not as a dependency. And reinforce the $K \to
K_{\text{eff}}$ distinction here, since it ties the whole methods section together.
\end{itemize}

\subsection*{Slide 22 --- Impact / key findings}
\scriptlabel Four findings, fast. \key{(1) Supervision beats the two-step} --- \num{0.81} vs
\num{0.72} synthetic, all five cohorts real. \key{(2) Preprocessing is decisive} --- per-platform
z-std or $\beta \to 0$; 10 of 12 alternatives collapse. \key{(3) Parsimony with signal} --- two
active programs vs $\sim$20, exact external scoring. \key{(4) Mean external C \num{0.636}}, fully
Bayesian, no $\alpha$. Then \key{honest limitations}: small noisy cohorts (Moffitt $\approx$\num{0.55});
modest real margin (EBMF strong within-training); \key{biology not yet annotated --- pathway
enrichment is next}.

\begin{itemize}
\item \sectionlabel{Main Takeaway:} Supervision helps clearly on synthetic and across all five real
cohorts; preprocessing is decisive; the model is parsimonious and fully Bayesian.
\item \sectionlabel{Emphasis:} Deliver the limitations honestly --- small-cohort noise, the modest
real-data margin, and the not-yet-done biological annotation. It strengthens credibility.
\item \sectionlabel{Good Line to Stress:} ``Two interpretable prognostic programs, exact external
scoring, no tuning parameter --- competitive with the field at a fraction of the complexity.''
\item \sectionlabel{Likely Question:} ``If the real-data margin is modest, why believe it?'' Because
the synthetic study isolates supervision's effect with known truth ($0.81$ vs $0.72$), and on real
data the projection model wins in \emph{all five} cohorts, not on average only --- consistency
across independent cohorts is the evidence, alongside parsimony and exact scoring.
\end{itemize}

\subsection*{Slide 23 --- Future directions \& summary}
\scriptlabel We \key{addressed two of April's next-steps}: CV-based $K$ selection and the
completed multi-cohort simulation; convergence is now on the full ELBO. \key{Next: pathway
enrichment} to name the two programs biologically, then \key{manuscript preparation}. The recommended
model in one line: \key{projection $(YF)\beta$, no cohort indicator, per-platform z-std, $K = 7$
($K_{\text{eff}} = 2$), mean external C \num{0.636}}.

\begin{itemize}
\item \sectionlabel{Main Takeaway:} The April next-steps are addressed; the remaining work is
biological annotation and the manuscript.
\item \sectionlabel{Emphasis:} Priors are point-exponential / normal --- not point-normal. Frame the
preprocessing point as ``crucial for performance and for avoiding precision domination,'' not jargon.
\item \sectionlabel{Good Line to Stress:} ``The next stage is turning two statistical programs into
interpretable biology.''
\end{itemize}

\subsection*{Slide 24 --- Discussion}
\scriptlabel Stop here, thank the group, and open the floor. If discussion turns technical, I have
appendix slides on the cohort-indicator contrast, on convergence and the survival-activity
magnitudes, and on the relationship to the lab's penalized survival-MF method.

\begin{itemize}
\item \sectionlabel{Main Takeaway:} Open for questions; note that deeper technical detail is in the
appendix.
\end{itemize}

\section*{Figure \& Table Analysis Guide}

This section decodes every figure and table in the deck with a consistent structure: what it
presents, how to interpret it, and the takeaway, followed by a key line to say and a likely
question. All numbers below come from the result CSVs and match the deck.

\subsection*{Figure 1 --- Factor recovery (synthetic) [Slide 15]}
\sectionlabel{What it presents.} A $2 \times 3$ grid of bar panels. Columns are the three scenarios
(all-shared, all-cohort-specific, shared \& cohort-specific); the top row is shared (prognostic)
factors and the bottom row is study-specific factors. The $y$-axis is factor recovery: the mean best
absolute correlation between each estimated program and the true program it matches, $0$ to $1$,
mean $\pm$ SE over five seeds. Bars are the five model arms (projection base/cohort, loadings
base/cohort, EBMF).

\sectionlabel{How to interpret it.} A tall bar means the estimated program closely reproduces a true
program. In the realistic hybrid scenario the projection model recovers shared factors at about
$0.89$ and study-specific factors at about $0.84$; the loadings model is lower ($\approx 0.72$ shared,
$\approx 0.67$ specific) and the unsupervised EBMF is lower still on shared factors ($\approx 0.68$).
In the all-shared scenario the projection model recovers shared factors at $\approx 0.92$.

\sectionlabel{The takeaway.} The supervised projection model reconstructs the true programs ---
especially the prognostic shared ones --- more faithfully than the loadings model or the
unsupervised baseline.

\begin{itemize}
\item \sectionlabel{Key line to say:} ``Recovery is just a correlation between each estimated program
and its true counterpart --- higher is better, and the projection model is highest on the shared
prognostic factors.''
\item \sectionlabel{Likely question:} ``Why is recovery not 1.0 even with supervision?'' Finite
sample size and noise; matrix factorization recovers programs up to small cross-loading, so $0.9$
with everything else near zero is clean recovery.
\end{itemize}

\subsection*{Figure 2 --- Specificity-classification accuracy (synthetic) [Slide 15]}
\sectionlabel{What it presents.} A $1 \times 3$ grid, one panel per scenario. The $y$-axis is
specificity-classification accuracy: the fraction of recovered factors correctly typed as shared
versus study-specific, judged from per-cohort loading energy alone (mean $\pm$ SE over five seeds).

\sectionlabel{How to interpret it.} This is a classification score, not a correlation. A factor that
loads across all cohorts should be called shared; one that loads in a single cohort should be called
specific. In the hybrid scenario the projection model reaches accuracy up to $1.00$ (cohort variant)
and $\approx 0.97$ (base); the loadings and EBMF arms are lower ($0.70$ to $0.83$).

\sectionlabel{The takeaway.} The projection model not only reconstructs programs but assigns them the
correct architectural type, which is why an explicit cohort indicator buys little.

\begin{itemize}
\item \sectionlabel{Key line to say:} ``Recovery asks whether the program is right; this asks whether
its shared-versus-specific label is right. We report both because you can get one without the other.''
\item \sectionlabel{Likely question:} ``Why do we need this in addition to recovery?'' A model could
reconstruct a program faithfully yet mistype it as shared when it is cohort-specific --- that would
be good programs with wrong structure. Both must hold.
\end{itemize}

\subsection*{Figure 3 --- Factor-recovery heatmap (synthetic) [Slide 16]}
\sectionlabel{What it presents.} Two side-by-side heatmaps of the patient-loading matrix $L$ in the
realistic hybrid scenario (seed 1). Rows are patients grouped by cohort (left annotation bar),
columns are factors $F1$--$F6$, with top annotations for true factor type (shared/prognostic vs.\
study-specific) and $|\hat\beta|$. Left: the true block-sparse $L$. Right: the estimated $L$ from the
projection model with no cohort indicator.

\sectionlabel{How to interpret it.} In a shared factor's column, both cohort blocks light up; in a
study-specific factor's column, only one block does. The estimated heatmap on the right reproduces
that block pattern from the true heatmap on the left.

\sectionlabel{The takeaway.} The model recovers the intended shared/specific block structure directly
from the data, without indicator columns --- the visual basis for Figures 1 and 2.

\begin{itemize}
\item \sectionlabel{Key line to say:} ``Shared programs are active in both cohort blocks, specific
programs in one --- and the estimate on the right matches the truth on the left.''
\end{itemize}

\subsection*{Figure 4 --- Held-out C-index by scenario (synthetic) [Slide 17]}
\sectionlabel{What it presents.} A $1 \times 3$ grid, one panel per scenario, of held-out C-index by
model arm (mean $\pm$ SE over five seeds) with a dashed line at chance ($0.50$). The unsupervised
EBMF $\to$ Cox baseline is shown alongside the supervised arms; values are labeled above the bars.

\sectionlabel{How to interpret it.} Higher is better; $0.50$ is chance. In the hybrid (realistic)
scenario: projection $0.81$, loadings $0.76$, EBMF $0.72$. In the all-cohort-specific negative
control, all arms sit at about $0.54$. In all-shared, projection is highest ($\approx 0.87$).

\sectionlabel{The takeaway.} Supervision improves discrimination in scenarios with real prognostic
structure and fabricates none in the null control.

\begin{itemize}
\item \sectionlabel{Key line to say:} ``$0.81$ for supervision versus $0.72$ for the two-step in the
realistic scenario --- and everyone at chance in the null, which is the safety check.''
\item \sectionlabel{Likely question:} ``Is the $0.81$-vs-$0.72$ gap meaningful given five seeds?''
The bars carry standard errors; the ordering (projection $>$ loadings $>$ EBMF) is consistent across
seeds, and the null-control collapse to chance confirms it is signal, not noise.
\end{itemize}

\subsection*{Figure 5 --- PDAC gene-weight heatmap (real) [Slide 18]}
\sectionlabel{What it presents.} A gene-weight heatmap of the top 30 genes for the recommended model
(projection $(YF)\beta$, per-platform z-standardization, survival-ranked gene selection over 2064
genes, no cohort indicator, $K = 7$). Columns are the seven programs; weights are non-negative
(point-exponential prior); survival-active programs ($|\hat\beta| > 10^{-3}$) are starred.

\sectionlabel{How to interpret it.} Each column is a gene program --- a coordinated set of
high-weight genes. Two columns are starred: Program 3 (protective) and Program 7 (adverse). The
largest expression program is a collagen/ECM stromal block and is \emph{not} starred --- it explains
expression but not survival.

\sectionlabel{The takeaway.} Supervision separates expression-only structure (the big stromal
program) from prognostic structure (the two starred programs) --- a separation a PCA-then-test
workflow cannot make.

\begin{itemize}
\item \sectionlabel{Key line to say:} ``The biggest expression program is not prognostic; only two
small programs carry survival signal. That separation is the payoff of supervision.''
\item \sectionlabel{Likely question:} ``Why are the weights all non-negative?'' Because $L$ and $F$
use a point-exponential prior, which constrains gene weights and loadings to be non-negative --- this
makes programs additive and interpretable.
\end{itemize}

\subsection*{Figure 6 --- Kaplan--Meier curves (real) [Slide 19]}
\sectionlabel{What it presents.} Kaplan--Meier survival curves for the two survival-active programs,
patients split at the median projection score $(YF)_k$ in the training cohorts. Each panel shows
high- vs.\ low-activation groups with a log-rank $p$-value and the empirical direction
(adverse/protective).

\sectionlabel{How to interpret it.} For the adverse program (Program 7, MET/ITGA3/glycolytic), the
high-activation group has worse survival, log-rank $p = 0.00047$. For the protective program
(Program 3, epithelial differentiation), the high-activation group has better survival, $p = 0.0077$.
The split is on the projection score, the quantity the risk model uses, not the raw loading.

\sectionlabel{The takeaway.} Both active programs produce clinically meaningful, statistically
significant survival stratification on the model's own score.

\begin{itemize}
\item \sectionlabel{Key line to say:} ``We split on $(YF)_k$, the score the model actually uses ---
the adverse program separates patients at $p = 0.0005$, the protective at $p = 0.008$.''
\item \sectionlabel{Likely question:} ``Why label by the curve rather than the $\beta$ sign?'' In the
joint model the coefficient sign for Program 7 can invert relative to its marginal direction because
the prognostic programs are correlated (a suppression effect). The empirical KM direction is the
honest, interpretable label, so we use it and stratify by the projection.
\end{itemize}

\subsection*{Figure 7 --- External C-index across five cohorts (real) [Slide 20]}
\sectionlabel{What it presents.} A grouped bar chart: cohort groups (Dijk, Moffitt, PACA-AU array,
PACA-AU seq, Puleo, and Mean), three bars per group --- loadings $L\beta$, projection $(YF)\beta$
(recommended), and unsupervised EBMF $\to$ Cox --- with a dashed chance line at $0.50$.

\sectionlabel{How to interpret it.} Per-cohort projection-model values: Dijk $0.657$, Moffitt
$0.550$, PACA-AU array $0.651$, PACA-AU seq $0.675$, Puleo $0.645$; mean $0.636$. The loadings model
means $0.622$; the unsupervised two-step means $0.564$ (per cohort $0.575$, $0.542$, $0.549$,
$0.562$, $0.591$). The projection model exceeds the two-step in every cohort. Moffitt is the weakest
($0.55$), reflecting a small, noisy microarray cohort.

\sectionlabel{The takeaway.} The real-data analog of Figure 4: the supervised projection model beats
the unsupervised two-step in all five held-out cohorts.

\begin{itemize}
\item \sectionlabel{Key line to say:} ``Mean $0.636$ for the projection model versus $0.564$ for the
two-step --- and it wins in all five cohorts, not just on average.''
\item \sectionlabel{Likely question:} ``Why is Moffitt so low?'' Small cohort, microarray, fewer
events --- per-cohort C-index is noisy at that sample size. It is still above chance, and the mean
and the all-cohort win are the robust signal.
\end{itemize}

\subsection*{Table --- Model comparison summary [Slide 21]}
\sectionlabel{What it presents.} One row per model --- projection (recommended), loadings, projection
+ cohort indicator, and unsupervised two-step --- with columns for linear predictor, preprocessing,
cohort indicator, $K$, $K_{\text{eff}}$, synthetic C (hybrid, 5 seeds), and real external C (mean
over 5 cohorts).

\sectionlabel{How to interpret it.} Read across the recommended (bold) row: projection $(YF)\beta$,
per-platform z-std + survival-ranked selection, no cohort indicator, $K = 7 \to K_{\text{eff}} = 2$,
synthetic $0.81$, real $0.636$. Compare down the last two columns: loadings $0.76$ / $0.622$;
projection + cohort indicator $0.81$ / $0.614$ (one extra factor, no gain); two-step $0.72$ / $0.564$.

\sectionlabel{The takeaway.} The recommended configuration wins on both the synthetic and the real
external axis, with the simplest survival structure and exact scoring.

\begin{itemize}
\item \sectionlabel{Key line to say:} ``Best on both axes, fewest active programs, exact external
scoring --- that is the recommendation.''
\end{itemize}

\subsection*{Appendix Figure A1 --- Cohort-indicator gene-weight heatmap [Slide A1]}
\sectionlabel{What it presents.} The gene-weight heatmap (top 40 genes) for the projection model with
the cohort-membership indicator added, $K = 8$, survival-active programs starred.

\sectionlabel{How to interpret it \& takeaway.} The survival-active structure is essentially
unchanged from the recommended (no-indicator) model --- the extra factor slot absorbs a platform
offset, not new biology. This confirms the offset is already removed by per-platform normalization.

\begin{itemize}
\item \sectionlabel{Key line to say:} ``Adding the indicator changes the bookkeeping, not the
biology --- the same two programs remain active.''
\end{itemize}

\subsection*{Appendix Figure A2 --- Cohort-indicator Kaplan--Meier [Slide A1]}
\sectionlabel{What it presents.} Kaplan--Meier stratification for the cohort-indicator model, split
by projection score, same format as Figure 6.

\sectionlabel{How to interpret it \& takeaway.} The survival separation is essentially identical to
the recommended model. Together with the external C-index ($0.614$ vs $0.636$), this shows the
indicator spends a factor without improving prognosis under per-platform z-standardization.

\subsection*{Appendix Figure A3 --- Training RMSE convergence traces [Slide A2]}
\sectionlabel{What it presents.} Reconstruction-RMSE traces across model configurations as a function
of CAVI iteration.

\sectionlabel{How to interpret it \& takeaway.} RMSE descends and stabilizes to a plateau; the actual
stopping rule is the relative ELBO change falling below $10^{-5}$ (with $\Delta L$, $\Delta\beta$
tracked as diagnostics). A small RMSE rise/plateau is expected because the joint objective trades a
little reconstruction for survival-relevant, regularized factors --- the ELBO is increasing
throughout.

\begin{itemize}
\item \sectionlabel{Likely question:} ``If RMSE plateaus or ticks up, is that a failure?'' No --- RMSE
is a by-product; the ELBO is the objective and it increases. The model trades a little expression fit
for survival-relevant structure.
\end{itemize}

\subsection*{Appendix Figure A4 --- Survival-activity magnitudes $|\hat\beta_k|$ [Slide A2]}
\sectionlabel{What it presents.} Bars of $|\hat\beta_k|$ per program, one facet per model, with red
for survival-active and gray for expression-only, and a dashed activity threshold.

\sectionlabel{How to interpret it \& takeaway.} Most programs fall below the $10^{-3}$ threshold
(expression-only); only a small number clear it (survival-active). This visualizes the
spike-and-slab shrinkage and how $K_{\text{eff}}$ is read off $|\hat\beta|$. For the recommended
model, Program 7 has the largest magnitude ($|\hat\beta| = 0.041$) and Program 3 the next
($|\hat\beta| = 0.011$).

\subsection*{Appendix Table --- Relation to the penalized survival-MF method (DeSurv) [Slide A3]}
\sectionlabel{What it presents.} A side-by-side comparison of the lab's penalized supervised-NMF+Cox
method and the present Bayesian model: estimation, what is tuned, shrinkage, uncertainty, and noise
model, with headline numbers.

\sectionlabel{How to interpret it \& takeaway.} Same design and same data --- both jointly fit
supervised NMF + Cox with non-negative programs, share the TCGA + CPTAC training set and the same
five external cohorts, score by projection, and select by CV + 1-SE. The differences: the penalized
method tunes $k$, $\alpha$, and $\lambda$ by Bayesian optimization, while the Bayesian model tunes
only $K$ and lets empirical-Bayes priors learn the shrinkage, with posterior uncertainty and a
per-gene precision $\tau$. The penalized method reports a pooled external HR per SD of $1.50$ (95\%
CI $1.31$--$1.72$); the Bayesian model reports mean external C $= 0.636$.

\begin{itemize}
\item \sectionlabel{Key line to say:} ``This is the Bayesian sibling of the penalized method --- same
idea, same data, but adaptive priors, only $K$ to tune, and posterior uncertainty.''
\item \sectionlabel{Honest caveat:} The two report different primary metrics (pooled HR/SD vs.\ mean
external C) and use slightly different normalization (within-subject rank $\to$ 1970 genes vs.\
per-platform z-std $\to$ 2064). A matched, same-protocol head-to-head is a planned next step --- the
penalized method is not re-run here, so do not claim superiority today.
\end{itemize}

\section*{Slide-by-Slide Review Guide (Quick Q\&A)}

Use this as a rapid self-test before the talk. For each prompt, you should be able to give the answer
spontaneously.

\begin{itemize}
\item \actioncue{Slide 4 --- Background:} \emph{Why model expression and survival jointly instead of two-step?}
--- Unsupervised axes chase the largest variance, often batch; a low-variance but prognostic program
gets missed. Supervision steers factors toward survival-relevant structure during estimation.

\item \actioncue{Slide 6 --- The model:} \emph{What makes a program prognostic?} --- Its survival coefficient
$\beta_k \neq 0$. If $\beta_k = 0$ the program explains expression only. The baseline hazard $h_0(t)$
is non-parametric and cancels in the partial likelihood.

\item \actioncue{Slide 7 --- Two linear predictors:} \emph{What is the difference between $L\beta$ and
$(YF)\beta$?} --- $L\beta$ is fully joint but needs an OLS re-projection to score new patients (an
approximation). $(YF)\beta$ scores new patients with the identical training formula --- exact,
leakage-free; it is the recommended model.

\item \actioncue{Slide 7 --- Priors:} \emph{What priors do you use?} --- Point-exponential (non-negative
spike-and-slab) on $L$ and $F$, normal on $\beta$, closed-form MLE for $\tau$. Not point-normal.

\item \actioncue{Slide 8 --- ELBO:} \emph{What is the ELBO used for?} --- It is the fitting objective CAVI
maximizes; convergence is its relative change below $10^{-5}$. It is \emph{not} the model-selection
criterion for $K$ --- that is cross-validation.

\item \actioncue{Slide 9 --- CAVI / convergence:} \emph{How do you know it converged?} --- Relative ELBO change
below $10^{-5}$ after a five-iteration burn-in; mean absolute changes in $L$ and $\beta$ (both below
$10^{-3}$) are diagnostic only.

\item \actioncue{Slide 10 --- Preprocessing:} \emph{Why is per-platform normalization essential?} --- The survival
precision $A_k = \sum_i w_i (YF)_{ik}^2$ is swamped by between-platform variance otherwise, collapsing
$\beta \to 0$. 10 of 12 non-per-platform pipelines collapse. We normalize each platform before
merging, leaving 2064 genes.

\item \actioncue{Slide 11 --- Selecting $K$:} \emph{How is $K$ chosen, and what is $K_{\text{eff}}$?} --- $K = 7$
by five-fold CV C-index, 1-SE rule, floor $K \ge 3$. $K_{\text{eff}} = 2$ is the subset with
$|\hat\beta| > 10^{-3}$.

\item \actioncue{Slide 12 --- Cohort indicator:} \emph{Why not use it?} --- Under per-platform z-std the offset is
already removed; it costs a factor ($K: 7 \to 8$) and lowers external C ($0.636 \to 0.614$).

\item \actioncue{Slide 15 --- Synthetic recovery:} \emph{What do the two synthetic recovery plots show?} ---
Recovery = correlation of estimated vs.\ true programs; specificity accuracy = correct shared/specific
typing. Projection is high on both.

\item \actioncue{Slide 17 --- Synthetic C-index:} \emph{What is the headline simulation result?} --- Hybrid
scenario: projection $0.81$ vs loadings $0.76$ vs two-step $0.72$; null control all at $\approx 0.54$.

\item \actioncue{Slide 18 --- Real programs:} \emph{How many programs, how many prognostic?} --- 7 retained
($K$), 2 survival-active ($K_{\text{eff}}$): Program 3 protective, Program 7 adverse. The largest
expression program (stromal/ECM) is not prognostic.

\item \actioncue{Slide 19 --- Real survival:} \emph{What are the survival results?} --- Split on median $(YF)_k$:
adverse program log-rank $p = 0.00047$, protective $p = 0.0077$; PH not violated.

\item \actioncue{Slide 20 --- External C-index:} \emph{What is the real-data headline?} --- Mean external C
$= 0.636$ (projection) vs $0.622$ (loadings) vs $0.564$ (two-step); projection wins in all five
cohorts.

\item \actioncue{Slides 22--23 --- Impact:} \emph{What are the limitations?} --- Small noisy external cohorts (Moffitt
$\approx 0.55$); modest real-data margin (EBMF strong within-training $\approx 0.63$); programs not
yet biologically annotated.
\end{itemize}

\section*{Anticipate These Audience Questions}

\begin{itemize}
\item \sectionlabel{``Why point-exponential on $L,F$ and normal on $\beta$ --- why not point-normal?''}
Programs and loadings should be non-negative and additive for interpretability, so a non-negative
point-exponential spike-and-slab is the natural prior on $L$ and $F$; it also gives sparse, readable
gene weights. The survival coefficients can take either sign, so a normal (Gaussian) prior is
appropriate there. Point-normal was a misstatement in the April materials --- this is the correction.

\item \sectionlabel{``How exactly is $K$ chosen, and how do you get $K_{\text{eff}}$?''} $K$ is chosen
by five-fold cross-validation: held-out C-index across a grid of $K$, the 1-SE rule (smallest $K$
within one SE of the best), and a biological floor of $K \ge 3$, giving $K = 7$. $K_{\text{eff}}$ is
then the number of those programs whose posterior $|\hat\beta_k|$ exceeds the activity threshold
$10^{-3}$ --- here 2. $K$ is reconstruction capacity; $K_{\text{eff}}$ is the prognostic subset.

\item \sectionlabel{``Why report two linear predictors --- isn't one enough?''} $L\beta$ is the
fully-joint reference where survival shapes the loadings most directly, but scoring a new patient
needs an OLS re-projection, an approximation. $(YF)\beta$ scores new patients with the identical
training formula --- exact and leakage-free --- so it is the recommended model. Reporting both
quantifies what we trade for exact scoring (very little: $0.636$ vs $0.622$ external).

\item \sectionlabel{``Is the ELBO how you choose the number of factors?''} No. The ELBO is the
\emph{fitting} objective --- CAVI maximizes it and convergence is when its relative change drops below
$10^{-5}$. The number of programs is chosen separately by cross-validated C-index. Conflating the two
was a point of confusion previously; keep them distinct.

\item \sectionlabel{``Why not include the cohort indicator --- doesn't it handle batch?''} It does, but
so does per-platform z-standardization, and they solve the same problem. Under per-platform z-std the
platform offset is already removed, so the indicator only spends a factor slot ($K: 7 \to 8$) and
slightly lowers external C ($0.636 \to 0.614$). We would only need it if we could not normalize per
platform.

\item \sectionlabel{``Have you tried other preprocessing --- does it really matter that much?''} Yes,
in an 18-configuration benchmark: any pipeline that does not z-standardize per platform before merging
(joint z-std, log-only, or joint quantile normalization without rank transform) collapses $\beta \to 0$
for both linear predictors --- 10 of 12 such configurations. The mechanism is that between-platform
variance dominates the survival precision $A_k = \sum_i w_i (YF)_{ik}^2$. Per-platform z-std before
merging is what makes the model work.

\item \sectionlabel{``How does this compare to the lab's penalized supervised-MF (DeSurv) method?''}
It is the Bayesian counterpart: same supervised NMF + Cox idea, same TCGA + CPTAC training set and
same five external cohorts, same projection-based scoring and CV + 1-SE selection, and we borrowed the
combined-rank gene-selection recipe. The differences are real: the penalized method tunes $k$,
$\alpha$, and $\lambda$ by Bayesian optimization; the Bayesian model tunes only $K$ and lets
empirical-Bayes priors learn the shrinkage, adds posterior uncertainty, and uses a per-gene noise
precision. They report different primary metrics (pooled HR/SD $= 1.50$ vs.\ mean external C $= 0.636$),
so a matched same-protocol head-to-head is a planned next step --- not a superiority claim today.

\item \sectionlabel{``Does the joint $\beta$ sign always match the survival direction you report?''}
Not necessarily. For the adverse program (Program 7), the joint-model coefficient sign can run
opposite to its marginal adverse direction --- a suppression effect among correlated prognostic
programs. We therefore label each program by the empirical survival direction the Kaplan--Meier curve
shows and stratify by the $(YF)$ projection, not the raw $\beta$ sign. This is deliberate, not an
inconsistency.

\item \sectionlabel{``Five external cohorts and a modest margin --- what would convince you it works?''}
Three things, in order of strength: the controlled simulation isolates supervision's effect with known
ground truth ($0.81$ vs $0.72$); the projection model wins in \emph{all five} independent real cohorts,
not on average only; and the recovered programs are parsimonious and biologically coherent
(MET/integrin/glycolytic adverse, epithelial protective). Pathway enrichment is the next step to
confirm the biology.

\item \sectionlabel{``How was the survival term made conjugate to the Bayesian updates?''} By a
second-order Taylor expansion of the Cox partial log-likelihood in the linear predictor $\eta = (YF)\beta$.
That produces a Gaussian-form working response and weight (the Cox working weight $w_i$), which slots
into the EBNM normal-means updates --- the same machinery as the genomics term.
\end{itemize}

\section*{Appendix Slide Scripts}

\subsection*{Slide A1 --- Appendix: Cohort-indicator contrast}
\scriptlabel If someone asks what the cohort indicator changes, this is the slide. Adding it under
per-platform z-standardization spends a factor slot --- $K$ goes from 7 to 8 --- without improving
external C-index ($0.636 \to 0.614$), and the biological survival-active structure is essentially
unchanged. That confirms the platform offset is already removed by normalization.

\begin{itemize}
\item \sectionlabel{Main Takeaway:} The indicator is redundant under per-platform normalization.
\end{itemize}

\subsection*{Slide A2 --- Appendix: Convergence \& survival-activity magnitudes}
\scriptlabel If the discussion turns to optimization behavior, these traces show RMSE stabilizing to
a plateau while the ELBO --- the actual objective --- increases; stopping is governed by the relative
ELBO change below $10^{-5}$. The $|\hat\beta_k|$ bars show the spike-and-slab shrinkage: most programs
fall below the activity threshold, only a couple clear it, which is how $K_{\text{eff}}$ is read off.

\begin{itemize}
\item \sectionlabel{Main Takeaway:} Convergence is judged on the ELBO; $K_{\text{eff}}$ is the few
programs with $|\hat\beta|$ above threshold.
\end{itemize}

\subsection*{Slide A3 --- Appendix: Relation to the penalized survival-MF method}
\scriptlabel This is the backup for the inevitable comparison question. Be collegial and precise: this
is the Bayesian sibling of the penalized supervised-NMF+Cox method --- same idea, same training data,
same five external cohorts, same projection scoring and CV + 1-SE selection, and we borrowed its
combined-rank gene selection. The genuine differences are adaptive empirical-Bayes priors (only $K$ to
tune, no $\alpha$ or $\lambda$), posterior uncertainty, and a per-gene noise model. The two report
different primary metrics, so a matched head-to-head on one protocol is the honest next step rather
than a superiority claim today.

\begin{itemize}
\item \sectionlabel{Main Takeaway:} Same problem and data, Bayesian estimation; differences are
adaptive priors, fewer tuned knobs, and uncertainty. Head-to-head is planned.
\end{itemize}

\end{document}
